\section{Business/societal Interpretation \& Value}

The analysis demonstrated that laundering activity was not distributed uniformly across
payment formats, originating banks and account relationships. These patterns illustrate
how distributed transaction analytics can support financial crime monitoring by
identifying transaction characteristics that may warrant further investigation.

\subsection{Value to Financial Institutions}

For financial institutions, the analysis demonstrates how large-scale transaction data
can be used to prioritise transactions and relationships for further investigation. For
example, payment formats with higher observed laundering rates can be incorporated
into risk-monitoring rules, while account relationships showing unusual frequency or
reciprocal activity can be subjected to additional review. Bank-level aggregation can
also assist institutions in identifying concentrations of labelled laundering activity within
their transaction-processing environment.

The key value of the distributed approach is scalability. Rather than analysing
transactions individually, PySpark can aggregate large numbers of records across
multiple dimensions, allowing AML analysts to examine patterns across payment
formats, institutions and account relationships.

\subsection{Value to Regulators and Financial-crime Researchers}

The approach can also support regulators and financial-crime researchers by providing
a framework for analysing transaction patterns at scale. Aggregated indicators can
assist with identifying areas requiring further investigation and can support the
development and evaluation of AML monitoring approaches.

\subsection{Societal value}

Effective AML monitoring can contribute to the integrity of the financial system by
improving the identification of potentially illicit transaction activity. More efficient
prioritisation of transactions for investigation may reduce the resources required to
review large volumes of transactions and allow investigators to focus on patterns that
warrant further examination. Improved identification of suspicious transaction activity
may also contribute to reducing the financial impact of fraud, organised crime and other
illicit activities that depend on the movement of funds through the financial system.

\subsection{Limitations}

Two important limitations affect the interpretation of the results.

\begin{enumerate}
    \item Synthetic data: The dataset is synthetic and therefore may not fully represent the behaviour, complexity or diversity of real-world financial transactions and money-laundering techniques.
    \item Historical Snapshot: The processed dataset covers a specific period between August 2022 and November 2022 and therefore may not reflect changing money-laundering strategies over time.
\end{enumerate}